% Carátula y primera página del Informe Final según el modelo del Anexo 5 del
% Reglamento de Proyecto Final, en una sola hoja (la primera del informe).
%  - Mitad superior: la carátula del modelo (logo, universidad, facultad,
%    carrera, trabajo, autor) con el formato del informe.
%  - Mitad inferior: la "primera página" del modelo (cátedra, docentes, tutor,
%    fecha de examen y el recuadro del tribunal) tal cual el modelo, en negro.
% La tapa de Scribus (tapa-scribus.pdf) queda solo para la encuadernación.

% Logo de la UTN, vectorial (mismas proporciones que el de la tapa de Scribus).
\newcommand{\logoutn}[1]{%
  \begin{tikzpicture}[x=#1/1095, y=#1/1095]
    \fill[m2navy] (35,1099) arc[start angle=180, end angle=360, radius=508]
      -- (866,1099) arc[start angle=0, end angle=-180, radius=323] -- cycle;
    \fill[m2navy] (35,35) arc[start angle=180, end angle=0, radius=508]
      -- (866,35) arc[start angle=0, end angle=180, radius=323] -- cycle;
    \fill[m2navy] (455,0) rectangle (632,1104);
    \fill[m2navy] (55,464) rectangle (1030,634);
  \end{tikzpicture}}

\newgeometry{top=15mm, bottom=15mm, left=25mm, right=25mm}
\begin{titlepage}
  \sffamily\setlength{\parskip}{1pt}
  % ---- carátula -------------------------------------------------------------------------
  \begin{center}
    \logoutn{16mm}\par\vspace{3mm}
    {\Large\bfseries\color{m2navy} Universidad Tecnológica Nacional\par}
    {\large\color{m2navy} Facultad Regional La Rioja\par}
    \vspace{2mm}
    {\large\bfseries\color{m2teal} Carrera de Ingeniería Electrónica\par}
    \vspace{1mm}
    {\color{m2amber}\rule{\textwidth}{1.4pt}\par}
  \end{center}
  {\small\itshape\color{m2grey} Trabajo Final de Grado:\par}
  \begin{center}
    {\LARGE\bfseries\color{m2navy}\parbox{0.7\textwidth}{\centering\titulotrabajo}\par}
    \vspace{2mm}{\large\color{m2teal}\subtitulotrabajo\par}
    \vspace{2mm}
    {\color{m2amber}\rule{\textwidth}{1.4pt}\par}
  \end{center}
  {\small\itshape\color{m2grey} Autor:\par}
  \begin{center}
    {\large\color{m2navy}\autor\par}
  \end{center}

  \vfill
  % ---- primera página: docentes y firmas, tal cual el modelo -------------------------------
  \color{black}\normalsize
  \hfill La Rioja, \fechaentrega\par
  \vspace{1mm}
  {\bfseries Cátedra de Proyecto Final\par}
  \vspace{2mm}
  Prof.\ Ing.\ Pedro Chahín Rearte\par
  JTP.\ Ing.\ Oscar Almonacid\par
  \vspace{3mm}
  \begin{tabular}{@{}p{45mm}@{}p{\dimexpr\textwidth-45mm\relax}@{}}
    Trabajo Final de Grado: & \titulotrabajo \\[2mm]
    Autor:                  & \autor \\[2mm]
  \end{tabular}\par
  Docente Tutor: \tutor\par
  \vspace{2mm}
  Fecha de Examen: \underline{\hspace{7mm}}-\underline{\hspace{7mm}}-\underline{\hspace{14mm}}\par
  \vspace{2mm}
  \begin{center}
    \setlength{\tabcolsep}{6pt}\renewcommand{\arraystretch}{1}%
    \newcommand{\fila}[1]{\multicolumn{3}{|p{\dimexpr0.87\textwidth+4\tabcolsep+2\arrayrulewidth\relax}|}{#1}}%
    \begin{tabular}{|p{0.29\textwidth}|p{0.29\textwidth}|p{0.29\textwidth}|}
      \hline
      \fila{\rule{0pt}{6mm}%
        \hspace{6mm}(Leyenda ``Aprobado'') \dotfill\hspace*{6mm}} \\[2mm]
      \fila{\hspace{6mm}EL EXAMEN DE DEFENSA DEL PRESENTE TRABAJO
        FINAL DE GRADO} \\[2mm]
      \fila{\hspace{6mm}HA MERECIDO LA CALIFICACIÓN DE
        \makebox[45mm]{\dotfill} PUNTOS} \\[2mm]
      \fila{\hspace{6mm}CONCEPTO: \dotfill\hspace*{6mm}} \\[2mm]
      \hline
      \fila{\rule{0pt}{4.5mm}Tribunal Examinador:} \\
      \hline
      \rule{0pt}{4.5mm}\small Presidente & \small Vocal 1° & \small Vocal 2° \\[20mm]
      \hline
    \end{tabular}
  \end{center}
\end{titlepage}
\restoregeometry
